\section{Convergent harmonic symplectic models}\label{sec:harmonic-models}

We first explain how the local equations and the metric use the same
coordinates. The qualitative quadraticity of the deformation problem
is known. The point needed here is control of the coordinate change
relative to the central harmonic norm.

\subsection{Conventions and the harmonic recursion}
Let $(X,\omega_0)$ be a Ricci-flat K3 surface and
\[
 E_0=\bigoplus_{i=1}^k F_i\otimes W_i
\]
a locally free slope-polystable bundle. The $F_i$ are pairwise
nonisomorphic stable bundles of common slope. Put $r_i=\rank F_i$
and $m_i=\dim W_i$, and choose a direct-sum Hermitian--Einstein
metric $h_0$. Write
\[
 d=\ddbar_{\End E_0},\quad V=H^1(X,\End E_0),\quad
 G=\prod_i\mathrm{GL}(W_i),\quad K=\prod_i\mathrm U(W_i).
\]
Ineffective factors are removed when taking effective quotients.
Since $H^1(X,\OO_X)=0$, all first-order deformations preserve the
holomorphic determinant. We use the $h_0,\omega_0$ harmonic norm.
Degrees are $c_1(F)\cdot[\omega]$, with no division by the volume
of $X$. The HE determinant connection varies with $\omega$.

Let $H_q$ and $G_q$ be the harmonic projections and Green operators.
At a fixed Sobolev order above the multiplication threshold, the
analytic implicit-function theorem solves
\begin{equation}\label{eq:harmonic-recursion}
 a(x)=x-d^*G_2(a(x)\wedge a(x)),\qquad
 F(x)=H_2(a(x)\wedge a(x))
\end{equation}
for small $x\in V$. Thus $a(x)=x+O(|x|^2)$, and $F=0$ is the
full convergent Kuranishi obstruction. The operators commute with
the parallel multiplicity matrices, so $a$ and $F$ are
$G$-equivariant wherever both sides are defined. This is a local
equivariance statement; a Euclidean ball is not $G$-invariant.

The surface hypothesis has a useful consequence. Since there are no
$(0,3)$-forms, $I-H_2=dd^*G_2$ on $(0,2)$-forms. Applying $d$ to
\eqref{eq:harmonic-recursion} gives the exact identity
\begin{equation}\label{eq:whole-curvature}
 da(x)+a(x)\wedge a(x)=F(x).
\end{equation}
It holds before imposing integrability. We will use it to regard the
obstruction itself as a moment map.

\subsection{A symplectic coordinate change with cubic correction}
Normalize the parallel holomorphic two-form $\sigma_X$ compatibly
with the harmonic norm, and define on $(0,1)$-forms
\[
 \mathcal B(\alpha,\beta)=
 \int_X\sigma_X\wedge\tr(\alpha\wedge\beta).
\]
In an adapted unitary coframe the antilinear operator
\begin{equation}\label{eq:quaternionic-J}
 J(A_1\bar\theta^1+A_2\bar\theta^2)
   =-A_2^*\bar\theta^1+A_1^*\bar\theta^2
\end{equation}
satisfies $J^2=-1$ and
$\mathcal B(\alpha,J\beta)=\langle\alpha,\beta\rangle_{L^2}$.
It preserves harmonic forms. Indeed the common-slope HE equation
makes the induced connection on $\End E_0$ anti-self-dual; the
K\"ahler identities, or integration by parts in the displayed
pairing followed by $J^2=-1$, give closedness and coclosedness of
$J\beta$ when $\beta$ is harmonic. This is the usual quaternionic
structure in the ASD moduli construction \cite{itoh}.

Let $B=a^*\mathcal B$ on $V$. It is closed and holomorphic.
Pairing curvature with a parallel endomorphism identifies its
complex moment with $F$, by \eqref{eq:whole-curvature} and Stokes'
theorem. For
$H_\xi(a)=\int_X\sigma_X\wedge\tr(\xi a\wedge a)$,
graded cyclicity of the trace gives
\[
 dH_\xi(a)(b)=\int_X\sigma_X\wedge
       \tr\bigl(\xi(b\wedge a+a\wedge b)\bigr)
       =\mathcal B([\xi,a],b).
\]
We use this convention for both the linear and nonlinear complex
moments.

The form $B_0$ is the Serre pairing on $V$. The correction
$D a_x-I$ takes values in $\im d^*$ and is $O(|x|)$. A harmonic
form is $\mathcal B$-orthogonal to $\im d^*$, by
\eqref{eq:quaternionic-J} and ordinary harmonic orthogonality.
The two terms linear in $Da_x-I$ therefore vanish, leaving
\begin{equation}\label{eq:symplectic-quadratic-error}
 B_x=B_0+O(|x|^2).
\end{equation}

\begin{proposition}\label{prop:symplectic-normal-form}
There is a $K$-equivariant holomorphic coordinate change on a
neighborhood of zero in $V$ such that
\begin{equation}\label{eq:cubic-darboux}
 \Psi^*B=B_0,\qquad
 \Psi(x)=x+O(|x|^3),\qquad D\Psi_x=I+O(|x|^2).
\end{equation}
Under the fixed Serre identification of the obstruction target with
the dual parallel algebra,
\[
 F\circ\Psi=Q,
\]
where $Q$ is the quadratic complex moment of the linear action on
$(V,B_0)$. The equality identifies the full convergent obstruction
ideals. Stable-factor ranks and multiplicities are arbitrary.
\end{proposition}

\begin{proof}
Write $\Delta B=B-B_0$. The holomorphic radial homotopy supplies
a one-form $\alpha$ with $d\alpha=\Delta B$. In coordinates,
\[
 \alpha_x(v)=\int_0^1 t\,\Delta B_{tx}(x,v)\,dt.
\]
Thus $\alpha=O(|x|^3)$, and it is $K$-invariant. On one smaller
ball all forms $B_s=B_0+s\Delta B$, $0\le s\le1$, are
nondegenerate. Solve $\iota_{V_s}B_s=-\alpha$. The resulting
holomorphic vector fields are invariant and vanish to order three,
uniformly in $s$. Their flows exist through time one on a common
smaller neighborhood. Cartan's formula gives
\[
 \frac{d}{ds}\Psi_s^*B_s
   =\Psi_s^*(\Delta B+d\iota_{V_s}B_s)=0.
\]
Their order of vanishing gives the two estimates in
\eqref{eq:cubic-darboux}. Holomorphic $K$-equivariance extends
locally to $G$-equivariance.

The two maps $F\circ\Psi$ and $Q$ are moment maps for the same
linear action and the same form $B_0$. Pairing their difference
with every parallel endomorphism has zero differential. Both
vanish at zero, so their difference is zero. The statement is an
equality of analytic maps, hence of their generated ideals.
\end{proof}

Paired harmonic orthonormal bases identify the representation with
a doubled Ext quiver, simultaneously normalizing its norm and
complex pairing. To track ranks, let $\eta_X$ be the harmonic
$(0,2)$-form with $\int_X\sigma_X\wedge\eta_X=1$. The harmonic
representative of a multiplicity-trace covector $(q_i)$ is
\[
 \bigoplus_i r_i^{-1}I_{F_i}\otimes q_i\eta_X.
\]
These factors belong to the obstruction-target identification.
The real moment is paired with the ordinary multiplicity trace.
For a nearby polarization parameter $\eta$ its target is
\begin{equation}\label{eq:general-real-level}
 \zeta_i(\eta)=2\pi\bigl(r_i\mu_{\omega_\eta}(E_0)
                          -\deg_{\omega_\eta}F_i\bigr),
 \qquad \sum_i m_i\zeta_i(\eta)=0.
\end{equation}
A determinant-preserving lift subtracts a rank-weighted common
scalar, which acts trivially on the deformation representation.

The DG-Lie formality theorem of \cite[Theorem~5.1]{bmm} already
applies to $E_0$, because $\Aut(E_0)$ is reductive. Proposition
\ref{prop:symplectic-normal-form} specifies convergent symplectic
coordinates and their derivative, which are the data required in
the next comparison. Equal reduced Hilbert polynomials are not
part of its hypotheses.

\subsection{The actual harmonic metric}
The full local HYM theorem \cite[Theorem~1.1 and Sections~3,5]{wuebben15}
applies to the convergent zero locus $F^{-1}(0)$. It gives nested
coefficient neighborhoods and an open classifying homeomorphism to
the actual fixed-determinant HYM moduli, smooth on smooth quotient
pieces and covering every sufficiently nearby solution. At a
balanced coefficient $x$, its comparison form is
\begin{equation}\label{eq:complex-orbit-minimum}
 g_{\rm fin}([b],[b])=
       \inf_{\xi\in\Lie G}\|b-C_x\xi\|^2,
 \qquad C_x\xi=[\xi,x].
\end{equation}
The finite orbit operator is $C_x^*C_x$; its positive eigenvalues
measure the nonzero infinitesimal orbit directions.
The actual harmonic metric lies between $1-C(|\eta|+|x|^2)$ and
$1+C(|\eta|+|x|^2)$ times this form. The comparison concerns
tangents to actual smooth quotient pieces, not all Zariski tangent
vectors at a singular point. Constants depend on the fixed central
background, its complementary elliptic gaps and bounded smooth
parameter data. They do not depend on ratios of positive eigenvalues
of $C_x^*C_x$ or on the nearby stabilizer type.

We record precisely how the coordinate change preserves this order.
Pull back the flat real K\"ahler form by $\Psi$. By
\eqref{eq:cubic-darboux} its error is $O(r^2)$ and its moment is
$\mu\circ\Psi=\mu+O(r^4)$, where $r=|x|$. These forms satisfy
the exact Hamiltonian identity. The finite balancing theorem
\cite[Theorem~2.1]{wuebben14}, applied on balls of radius comparable
to $r$, changes the real orbit by $e^S$, $\|S\|=O(r^2)$, to put
the pulled-back and flat moments at the same central target.
The theorem includes the inverse correspondence and permits the
invariant analytic subset $Q^{-1}(0)$.

For every $\xi$, equivariance gives
$D\Psi_x C_x\xi=C_{\Psi(x)}\xi$ and
\[
 (1-Cr^2)\|b-C_x\xi\|^2
 \le\|D\Psi_x(b-C_x\xi)\|^2
 \le(1+Cr^2)\|b-C_x\xi\|^2.
\]
Infimizing preserves the inequalities. The derivative of the
varying balancing gauge contributes a complex orbit vector,
which disappears in this same infimum; its value has norm
distortion $1+O(r^2)$. Thus a derivative estimate for that gauge,
which need not be uniform at a stabilizer change, is unnecessary.

\begin{corollary}\label{cor:harmonic-comparison}
Put $\Q_\zeta=\{Q=0,\mu=\zeta\}/K$, with its flat quotient
metric. On fixed nested neighborhoods there are classifying
identifications $\Phi_\eta$ with the actual HYM moduli satisfying
\begin{equation}\label{eq:actual-quotient-comparison}
 (1-C(|\eta|+r^2))g^{\rm quot}_{\zeta(\eta)}
 \le\Phi_\eta^*g_\eta
 \le(1+C(|\eta|+r^2))g^{\rm quot}_{\zeta(\eta)}
\end{equation}
on every smooth quotient piece. Their images cover every sufficiently
nearby solution. The constants are independent of positive finite
scale ratios. The fixed holomorphic family
$x\mapsto(E_0,\ddbar_0+a(\Psi(x)))$ determines the classifying
maps; the HE gauge corrections preserve its isomorphism classes.
\end{corollary}

On smooth complex quotient charts the classifying maps are
holomorphic. Their open injectivity makes them biholomorphisms
onto their images. Statements about entire exceptional fibres
will require the additional compactness arguments below; they
are not consequences of a germ identification alone.

\subsection{Mukai and harmonic normalizations}
For a universal family $\cE$ on $X\times M$, fix
\begin{equation}\label{eq:mukai-marking}
 \theta_v(x)=\left[p_{M*}\bigl(\operatorname{ch}(\cE)
       \sqrt{\operatorname{td}_X}\,p_X^*(x^\vee)\bigr)\right]_2,
 \qquad x^\vee=(r,-c,s),\quad x\in v^\perp.
\end{equation}
The perpendicular is taken in the full integral Mukai lattice.
Twisting $\cE$ by a line from $M$ does not change this map on
$v^\perp$. We also use its real extension and the
Beauville--Bogomolov form $q$, including the rational dual form
on curve functionals.

The unnormalized real Dolbeault metric is half the usual real
connection $L^2$ metric. In this convention its K\"ahler class is
\begin{equation}\label{eq:harmonic-class}
 [\Omega_\eta]=2\pi^2\theta_v\left(
  0,[\omega_\eta],
       \frac{c_1(v)\cdot[\omega_\eta]}{\rank v}\right).
\end{equation}
To verify the factor, set
$\Delta(\cE)=c_2(\cE)-(r-1)c_1(\cE)^2/(2r)$.
Biswas--Schumacher's fixed-determinant formula
\cite[Proposition~4.3, (4.6)]{biswasschumacher} is
\[
 \frac{\omega^{\rm WP}}{2\pi^2}
 =-\frac1r p_{M*}\bigl(\operatorname{ch}_2(\End\cE)\omega_\eta\bigr)
 =2p_{M*}\bigl(\Delta(\cE)\omega_\eta\bigr).
\]
Their form is $iG_{j\bar k}\,dz_j\wedge d\bar z_k$, whereas
the real metric $\operatorname{Re}G$ has K\"ahler form
$\omega^{\rm WP}/2$. Expanding
\eqref{eq:mukai-marking} identifies the last pushforward with
the Mukai expression in \eqref{eq:harmonic-class}. This also
shows that base-line twists cancel. In particular cohomological
periods and quotient moment scales can now be compared without
an unspecified normalization constant.

